% 第四部：論證的前提（第貳部分，1359–1369）。出處：細目次第三節；W06_1019_導讀問答.md「第一項假設」一則。
\chapter{論證的前提：四項假設}

\qaentry{第四部之一}{四項假設的內容}

\begin{qaquestion}
本文的論證以哪四項假設為前提？作者為何要把論證寫成附條件的形式？
\end{qaquestion}

第貳部分前言列出四項假設：

\begin{quote}
We are to imagine a society with (1) democratic institutions in reasonably good working order, including a representative legislature elected on the basis of universal adult suffrage; (2) a set of judicial institutions, again in reasonably good order, set up on a nonrepresentative basis to hear individual lawsuits, settle disputes, and uphold the rule of law; (3) a commitment on the part of most members of the society and most of its officials to the idea of individual and minority rights; and (4) persisting, substantial, and good faith disagreement about rights (i.e., about what the commitment to rights actually amounts to and what its implications are) \ldots\pg{1360}
\end{quote}

接著寫：「I shall argue that, relative to these assumptions, the society in question ought to settle the disagreements about rights that its members have using its legislative institutions.」\pg{1360}摘要第 2 段已說明，論據並非絕對或無條件，以社會有運作良好的民主制度、多數公民認真對待權利為條件\pg{1346}；任一條件不成立，論證可能不成立，留待第柒部分處理。

\begin{longtable}{@{}L{0.16\linewidth}L{0.35\linewidth}L{0.22\linewidth}L{0.16\linewidth}@{}}
\toprule
假設 & 內容與限定 & 後文取用處 & 不成立時 \\
\midrule
\endhead
一、民主制度（1361--1362） & 成年人普選、定期公平選舉、大型審議機關、兩院與委員會、政黨；運作良好而未必完美；只指程序，不預設立法內容大體正義 & 第伍部分的立法機關答覆；第柒部分 & 1402--1403 \\
二、司法制度（1363--1364） & 獨立而運作良好的法院，已有權審查行政行為；非民選、非代議；回應具體訴訟、對抗式程序、遵循先例、合議庭以多數決作成決定 & 第肆部分 A 節；第伍部分的法院答覆 & 1402 \\
三、對權利的承諾（1364--1366） & 多數成員與公職人員認真對待個人與少數者的權利；權利懷疑論屬例外；已有一部權利法案，是否入憲、是否剛性不問 & 第陸、柒部分 & 1398；1404 \\
四、關於權利的分歧（1366--1369） & 就有哪些權利、其內容為何存在實質而善意的分歧，及於權利的核心適用 & 第參部分；第肆部分 B 節 & 1404 \\
\bottomrule
\end{longtable}

\qaentry{第四部之二}{第一、二項假設的平衡}

\begin{qaquestion}
第一項假設是否把立法程序理想化？作者如何避免乞題？
\end{qaquestion}

第一項假設限定其程度：

\begin{quote}
In general, I am assuming that the democratic institutions are in reasonably good order. They may not be perfect and there are probably ongoing debates as to how they might be improved.\pg{1361}
\end{quote}

同節末段把立法的內容排除在假設之外：作者不假定合理民主的立法機關所制定的法律大體良好或正義，他對立法與選舉安排運作良好的說法只涉及程序價值，不涉及結果價值\pg{1362}。提出第一項假設時，作者即以第二項假設平衡，設定法院同樣運作良好，表明不願乞題\pg{1362}。第肆部分前言另自承：對立法機關的樂觀描述是刻意的，用以對應法院支持者拒絕對法院作實際描述時的樂觀描述\pg{1379}。

\qaentry{第四部之三}{權利分歧與權利懷疑論}

\begin{qaquestion}
第四項假設所說的分歧，與道德相對主義或權利懷疑論有何不同？權利法案能否解決這種分歧？
\end{qaquestion}

作者所舉的分歧議題有墮胎、平權措施、重分配、犯罪嫌疑人的權利、宗教寬容、少數者的文化權利、選舉言論與支出的管制\pg{1367}。這些分歧出於善意，作者引 Rawls《政治自由主義》所稱的判斷的負擔\pg{1368}；承認分歧與道德相對主義無關\pg{1368}。分歧的雙方都認真對待權利，由第三項假設保證。權利法案與爭點有關，但不解決爭點；設有違憲審查時，分歧被構作成條文的解釋問題\pg{1368--1369}，第肆部分 B 節由此展開。

\begin{readernote}
第三、四兩項假設合起來，把「多數不尊重少數權利」的情形排在核心情形之外，留下的只有「雙方都尊重權利、對內容意見不同」的情形。違憲審查的辯護者最常舉的例子多半屬於前一種，聽者若覺得本文避開了真正的難題，可以指向這兩項假設與第柒部分。
\end{readernote}
